% Theoretical Foundations

\section{Theoretical Foundations}

\subsection{Fuzzy Set Theory for Molecular Evidence}

Let $U$ denote the universal set of molecular evidence features. A fuzzy set $A$ in $U$ is characterized by membership function $\mu_A: U \rightarrow [0,1]$ where $\mu_A(x)$ represents the degree of membership of element $x$ in set $A$ \cite{zadeh1965fuzzy}.

For molecular evidence vector $\mathbf{e} = (e_1, e_2, \ldots, e_n)^T$ where $e_i$ represents individual evidence components, membership functions are defined as:

\begin{equation}
\mu_{LOW}(e_i) = \begin{cases}
1 & \text{if } e_i \leq a \\
\frac{b - e_i}{b - a} & \text{if } a < e_i < b \\
0 & \text{if } e_i \geq b
\end{cases}
\end{equation}

\begin{equation}
\mu_{MEDIUM}(e_i) = \begin{cases}
0 & \text{if } e_i \leq c \text{ or } e_i \geq f \\
\frac{e_i - c}{d - c} & \text{if } c < e_i < d \\
1 & \text{if } d \leq e_i \leq e \\
\frac{f - e_i}{f - e} & \text{if } e < e_i < f
\end{cases}
\end{equation}

\begin{equation}
\mu_{HIGH}(e_i) = \begin{cases}
0 & \text{if } e_i \leq g \\
\frac{e_i - g}{h - g} & \text{if } g < e_i < h \\
1 & \text{if } e_i \geq h
\end{cases}
\end{equation}

where parameters $a, b, c, d, e, f, g, h$ define membership function boundaries with constraints $a < b \leq c < d \leq e < f \leq g < h$.

\subsection{Bayesian Network Structure}

Bayesian network $G = (V, E)$ consists of directed acyclic graph with vertex set $V$ representing random variables and edge set $E$ representing conditional dependencies \cite{pearl1988probabilistic}. For molecular identification, variables include:

$M$: molecular identity hypothesis
$S$: spectral evidence
$T$: structural evidence  
$P$: pathway evidence
$C$: confidence assessment

Joint probability distribution factorizes as:
\begin{equation}
P(M, S, T, P, C) = P(M) \cdot P(S|M) \cdot P(T|M) \cdot P(P|M) \cdot P(C|S, T, P)
\end{equation}

\subsection{Hybrid Fuzzy-Bayesian Inference}

Evidence variables $S, T, P$ are fuzzified using membership functions \cite{mamdani1975experiment}. For spectral evidence $S$ with value $s$, fuzzy representation becomes:
\begin{equation}
\tilde{S} = \{\mu_{LOW}(s)/LOW, \mu_{MEDIUM}(s)/MEDIUM, \mu_{HIGH}(s)/HIGH\}
\end{equation}

Fuzzy-Bayesian posterior calculation:
\begin{equation}
P(M = m_j | \tilde{E}) = \frac{\sum_{i} w_i \cdot P(E_i | M = m_j) \cdot P(M = m_j)}{\sum_{k} \sum_{i} w_i \cdot P(E_i | M = m_k) \cdot P(M = m_k)}
\end{equation}

where $\tilde{E} = \{\mu_i/E_i\}$ is fuzzy evidence set, $w_i = \mu_i$ are membership weights, and summation over $i$ covers all fuzzy evidence categories.

\subsection{Temporal Evidence Decay}

Evidence reliability decreases exponentially with time \cite{norris1997markov}:
\begin{equation}
R(t) = R_0 \cdot e^{-\lambda t}
\end{equation}

where:
- $R_0$: initial reliability at $t = 0$
- $\lambda$: decay constant (s$^{-1}$)
- $t$: time elapsed since evidence collection (s)

For 30-day half-life: $\lambda = \frac{\ln(2)}{30 \times 24 \times 3600} = 2.68 \times 10^{-7}$ s$^{-1}$

Time-adjusted posterior probability:
\begin{equation}
P(M = m_j | \tilde{E}, t) = \frac{\sum_{i} R_i(t) \cdot w_i \cdot P(E_i | M = m_j) \cdot P(M = m_j)}{\sum_{k} \sum_{i} R_i(t) \cdot w_i \cdot P(E_i | M = m_k) \cdot P(M = m_k)}
\end{equation}

\subsection{Uncertainty Propagation}

Total uncertainty $U_{total}$ combines aleatory uncertainty $U_{aleatory}$ and epistemic uncertainty $U_{epistemic}$ \cite{helton2004alternative}:
\begin{equation}
U_{total}^2 = U_{aleatory}^2 + U_{epistemic}^2
\end{equation}

Aleatory uncertainty from measurement noise:
\begin{equation}
U_{aleatory} = \sqrt{\sum_{i=1}^n \sigma_i^2}
\end{equation}

where $\sigma_i$ is measurement standard deviation for evidence component $i$.

Epistemic uncertainty from model limitations:
\begin{equation}
U_{epistemic} = \sqrt{\text{Var}[P(M|\tilde{E})]}
\end{equation}

Confidence interval for molecular identification \cite{efron1993bootstrap}:
\begin{equation}
CI = P(M = m_j | \tilde{E}) \pm z_{\alpha/2} \cdot U_{total}
\end{equation}

where $z_{\alpha/2}$ is critical value for confidence level $1-\alpha$.
